% GENERATED FILE. Edit canonical lessons, then run npm run generate:books.
% canonical-source-hash: fnv1a64:a1033dfa5dc6844a
% canonical-lessons: GE-C274-geruecht, GE-C274-quelle, GE-C274-behaupten, GE-C274-behauptung, GE-C274-bestaetigen

\chapter{Ein Gerücht}
\label{ch:ge-ein-gerucht}
\hlchaptermodality{274}

\begin{chapteropening}
\textbf{By the end of this chapter:} \emph{I can talk about how true a story is: a rumour, a source, to claim, a claim and to confirm.}

Complete the last lesson of chapter 274: I can talk about how true a story is: a rumour, a source, to claim, a claim and to confirm.
\end{chapteropening}

\section[das Gerücht]{das Gerücht — the rumour}

\label{lesson:GE-C274-geruecht}



\begin{quote}
Before the new one: say the German for to announce, then the German for to make public.
\end{quote}

\subsection*{You'll want to know: das Gerücht}
\textbf{das Gerücht} — \textquotedblleft{}the rumour\textquotedblright{}.

\begin{quote}
Im Dorf gab es ein Gerücht über den neuen Nachbarn.
\end{quote}

\begin{quote}
\emph{In the village there was a rumour about the new neighbour.}
\end{quote}

Plural: \emph{die Gerüchte}.

\begin{cousinweb}
\emph{Gerücht} is usually traced to Middle Low German \emph{gerüchte} \textquotedblleft{}outcry, talk\textquotedblright{}, built on the root of \emph{rufen} \textquotedblleft{}to call\textquotedblright{}, from Proto-Germanic \emph{hrōpaną} (EtymDB).
\end{cousinweb}

\begin{grammarlens}[title={What you've built}]
One more word for telling, explaining and describing at B1.
\end{grammarlens}

\subsection*{Your turn}
\begin{itemize}
\raggedright
  \item \emph{Say it:} \emph{das Gerücht}
  \item \emph{Say it:} \emph{das Gerücht}, once more
  \item \emph{Say it:} a sentence of your own with \emph{das Gerücht}
  \item \emph{Recall:} say the German for to announce, then the German for to make public, then say \emph{das Gerücht} again
\end{itemize}

\begin{tcolorbox}[breakable,colback=teal!4,colframe=teal!35!black,title={Before you move on}]
What does das Gerücht mean? (\textquotedblleft{}The rumour\textquotedblright{}.) Say it once more.
\end{tcolorbox}
\section[die Quelle]{die Quelle — the source; the spring}

\label{lesson:GE-C274-quelle}



\begin{quote}
Before the new one: say the German for to make public, then the German for the rumour.
\end{quote}

\subsection*{You'll want to know: die Quelle}
\textbf{die Quelle} — \textquotedblleft{}the source; the spring\textquotedblright{}.

\begin{quote}
Der Journalist nannte seine Quelle nicht.
\end{quote}

\begin{quote}
\emph{The journalist did not name his source.}
\end{quote}

Plural: \emph{die Quellen}. The first sense is a spring of water.

\begin{cousinweb}
\emph{Quelle} is inherited from Old High German \emph{quella}, which EtymDB traces to Proto-Germanic \emph{kwellōn} \textquotedblleft{}well, spring\textquotedblright{}: where water, and later news, comes from.
\end{cousinweb}

\begin{grammarlens}[title={What you've built}]
One more word for telling, explaining and describing at B1.
\end{grammarlens}

\subsection*{Your turn}
\begin{itemize}
\raggedright
  \item \emph{Say it:} \emph{die Quelle}
  \item \emph{Say it:} \emph{die Quelle}, once more
  \item \emph{Say it:} a sentence of your own with \emph{die Quelle}
  \item \emph{Recall:} say the German for to make public, then the German for the rumour, then say \emph{die Quelle} again
\end{itemize}

\begin{tcolorbox}[breakable,colback=teal!4,colframe=teal!35!black,title={Before you move on}]
What does die Quelle mean? (\textquotedblleft{}The source; the spring\textquotedblright{}.) Say it once more.
\end{tcolorbox}
\section[behaupten]{behaupten — to claim, to assert}

\label{lesson:GE-C274-behaupten}



\begin{quote}
Before the new one: say the German for the rumour, then the German for the source.
\end{quote}

\subsection*{You'll want to know: behaupten}
\textbf{behaupten} — \textquotedblleft{}to claim, to assert\textquotedblright{}.

\begin{quote}
Er behauptete, dass er nichts gesehen hatte.
\end{quote}

\begin{quote}
\emph{He claimed that he had seen nothing.}
\end{quote}

Often with \emph{dass} and the verb at the end. Reflexive, \emph{sich behaupten} means to hold one's own.

\begin{cousinweb}
\emph{behaupten} is usually traced to Middle High German \emph{behoubeten} \textquotedblleft{}to hold one's ground\textquotedblright{}, built on \emph{Haupt} \textquotedblleft{}head\textquotedblright{}, a cousin of English \emph{head} (EtymDB).
\end{cousinweb}

\begin{grammarlens}[title={What you've built}]
One more word for telling, explaining and describing at B1.
\end{grammarlens}

\subsection*{Your turn}
\begin{itemize}
\raggedright
  \item \emph{Say it:} \emph{behaupten}
  \item \emph{Say it:} \emph{behaupten}, once more
  \item \emph{Say it:} a sentence of your own with \emph{behaupten}
  \item \emph{Recall:} say the German for the rumour, then the German for the source, then say \emph{behaupten} again
\end{itemize}

\begin{tcolorbox}[breakable,colback=teal!4,colframe=teal!35!black,title={Before you move on}]
What does behaupten mean? (\textquotedblleft{}To claim, to assert\textquotedblright{}.) Say it once more.
\end{tcolorbox}
\section[die Behauptung]{die Behauptung — the claim, the assertion}

\label{lesson:GE-C274-behauptung}



\begin{quote}
Before the new one: say the German for the source, then the German for to claim.
\end{quote}

\subsection*{You'll want to know: die Behauptung}
\textbf{die Behauptung} — \textquotedblleft{}the claim, the assertion\textquotedblright{}.

\begin{quote}
Diese Behauptung ist falsch.
\end{quote}

\begin{quote}
\emph{This claim is false.}
\end{quote}

Plural: \emph{die Behauptungen}.

\begin{cousinweb}
\emph{Behauptung} is built on \emph{behaupten} \textquotedblleft{}to claim\textquotedblright{}, which rests on \emph{Haupt} \textquotedblleft{}head\textquotedblright{}, an inherited Germanic word (EtymDB).
\end{cousinweb}

\begin{grammarlens}[title={What you've built}]
One more word for telling, explaining and describing at B1.
\end{grammarlens}

\subsection*{Your turn}
\begin{itemize}
\raggedright
  \item \emph{Say it:} \emph{die Behauptung}
  \item \emph{Say it:} \emph{die Behauptung}, once more
  \item \emph{Say it:} a sentence of your own with \emph{die Behauptung}
  \item \emph{Recall:} say the German for the source, then the German for to claim, then say \emph{die Behauptung} again
\end{itemize}

\begin{tcolorbox}[breakable,colback=teal!4,colframe=teal!35!black,title={Before you move on}]
What does die Behauptung mean? (\textquotedblleft{}The claim, the assertion\textquotedblright{}.) Say it once more.
\end{tcolorbox}
\section[bestätigen]{bestätigen — to confirm}

\label{lesson:GE-C274-bestaetigen}



\begin{quote}
Before the new one: say the German for to claim, then the German for the claim.
\end{quote}

\subsection*{You'll want to know: bestätigen}
\textbf{bestätigen} — \textquotedblleft{}to confirm\textquotedblright{}.

\begin{quote}
Die Polizei bestätigte die Meldung.
\end{quote}

\begin{quote}
\emph{The police confirmed the news item.}
\end{quote}

Perfekt: \emph{hat bestätigt}. Also of a booking: \emph{Das Hotel bestätigte die Reservierung}.

\begin{cousinweb}
\emph{bestätigen} is usually traced to Middle High German \emph{bestætigen} \textquotedblleft{}to make firm\textquotedblright{}, built on \emph{stætec} \textquotedblleft{}steady\textquotedblright{}, from the root of \emph{stehen} \textquotedblleft{}to stand\textquotedblright{}; EtymDB lists no parent.
\end{cousinweb}

\begin{grammarlens}[title={What you've built}]
One more word for telling, explaining and describing at B1.
\end{grammarlens}

\subsection*{Your turn}
\begin{itemize}
\raggedright
  \item \emph{Say it:} \emph{bestätigen}
  \item \emph{Say it:} \emph{bestätigen}, once more
  \item \emph{Say it:} a sentence of your own with \emph{bestätigen}
  \item \emph{Recall:} say the German for to claim, then the German for the claim, then say \emph{bestätigen} again
\end{itemize}

\begin{tcolorbox}[breakable,colback=teal!4,colframe=teal!35!black,title={Before you move on}]
What does bestätigen mean? (\textquotedblleft{}To confirm\textquotedblright{}.) Say it once more.
\end{tcolorbox}
